\chapter{无锁正确性证明、回收与实现边界}

\section{为什么原理之后还要证明}

第 12 章已经说明了 SPSC、MPSC、MPMC、ABA、hazard pointer 和 epoch 的核心问题。但无锁结构只讲直觉仍然不够。真正能进入工程的无锁实现，必须给出三个证明面：线性化证明、进展证明和内存回收证明。缺少任意一面，结构都可能在小测试中正确，在压力、取消、线程暂停或对象池复用时失败。

\begin{lstlisting}[numbers=none,caption={无锁实现的三份证明}]
correctness:
  linearizability
  memory_order_happens_before
  reclamation_safety

progress:
  blocking / obstruction-free / lock-free / wait-free
  allocator_dependency
  backoff_and_starvation

engineering:
  stress_test
  sanitizer
  model_or_litmus_test
  metrics
  failure_replay
\end{lstlisting}

\section{Linearizability：每个操作在哪一瞬间生效}

线性化证明的任务，是为每个公开操作找一个瞬间，使并发历史等价于某个串行历史。对栈而言，push 成功 CAS 通常是线性化点；pop 成功 CAS 通常是线性化点；pop 发现空的线性化点是读取到空 head 的那个观察点，但前提是没有并发 push 已经在线性化。失败返回同样需要证明。

\begin{center}
\small
\begin{tabular}{p{0.22\linewidth}p{0.34\linewidth}p{0.32\linewidth}}
\toprule
操作 & 候选线性化点 & 证明要点 \\
\midrule
SPSC push 成功 & release 写 tail & slot 数据先于 tail 可见 \\
SPSC pop 成功 & release 写 head 或 acquire 观察 tail & 数据读取属于已发布区间 \\
Treiber push 成功 & CAS head 成功 & 新节点 next 指向旧 head \\
Treiber pop 成功 & CAS head 成功 & 被弹节点从结构中唯一移除 \\
MPMC enqueue & slot sequence 发布 Ready & 预留和可读提交要分开 \\
MPMC dequeue & 消费者 CAS 或 sequence 更新 & 元素只被一个消费者取得 \\
\bottomrule
\end{tabular}
\end{center}

证明不能停在“CAS 成功”。CAS 成功只是修改了某个原子对象。你还要证明这个修改与普通数据写入之间有 release/acquire 关系，证明没有另一个线程能同时取得同一元素，证明失败路径没有把未发布数据当成可见事实。

\section{内存序证明：发布的是普通数据，不是心情}

无锁结构通常用 relaxed 读写索引，用 release/acquire 发布数据。关键问题是：哪个普通写必须先于哪个普通读可见？SPSC push 中，生产者先写槽位，再 release 存储 tail；消费者 acquire 读取 tail 后读槽位。这样 slot 写入 happens-before 消费者读取。

\begin{lstlisting}[numbers=none,caption={SPSC push/pop 的 happens-before 证明}]
producer:
  buffer[tail] = value
  tail.store(next, release)

consumer:
  observed_tail = tail.load(acquire)
  value = buffer[head]

proof:
  release tail synchronizes-with acquire tail
  buffer write is sequenced-before release
  buffer read is sequenced-after acquire
  therefore consumer sees published value
\end{lstlisting}

若把 tail 存储放到写槽位之前，或者把消费者 acquire 改成没有同步关系的普通读，就可能读到未初始化或旧值。内存序不是装饰参数，而是把普通内存和原子状态连接起来的合同。

\section{ABA：地址相同不代表对象相同}

ABA 的本质是历史丢失。线程 T1 观察到 head 是地址 A；T2 把 A 弹出并释放；allocator 又把同一地址给了新节点；T1 再看 head 仍是 A，于是 CAS 成功。CAS 只比较位模式，不比较对象生命周期。解决 ABA 可以用 tagged pointer、hazard pointer、epoch reclamation、引用计数或固定数组 sequence。每种方案都转移成本。

\begin{center}
\small
\begin{tabular}{p{0.22\linewidth}p{0.34\linewidth}p{0.32\linewidth}}
\toprule
方案 & 解决方式 & 代价 \\
\midrule
tagged pointer & 指针加版本号 & 位宽、溢出、平台约束 \\
hazard pointer & 读者声明正在保护的地址 & 每线程槽位、扫描成本、保护顺序严格 \\
epoch reclamation & 旧 epoch 无活动读者后回收 & 暂停线程会阻塞回收 \\
固定数组 sequence & 地址不释放，用轮次区分状态 & 有界容量、槽位状态更复杂 \\
引用计数 & 对象存活由计数控制 & 原子成本和循环引用问题 \\
\bottomrule
\end{tabular}
\end{center}

\section{Hazard pointer：保护窗口必须重新读取}

Hazard pointer 的关键顺序是：先读取候选指针，发布到 hazard 槽位，再重新读取共享指针确认没有变化。少了重新读取，就可能保护了已经不再是结构头的对象；发布 hazard 太晚，就可能在发布前对象已经被回收。

\begin{lstlisting}[numbers=none,caption={hazard pointer 的最小保护顺序}]
retry:
  p = head.load(acquire)
  hazard_slot.store(p, release)
  if (head.load(acquire) != p):
      goto retry
  // p is protected while hazard_slot == p
  next = p->next
  if (CAS(head, p, next)):
      hazard_slot.store(nullptr, release)
      retire(p)
\end{lstlisting}

回收线程扫描所有 hazard 槽位。retired list 中不被任何 hazard 保护的节点才能真正 delete。这个算法不是免费：线程越多，扫描越贵；retired list 太小会频繁扫描，太大又占内存；线程退出时必须清理 hazard 槽位；异常和取消路径必须释放保护。

\section{Epoch reclamation：快，但怕暂停线程}

Epoch based reclamation 把时间分成 epoch。线程进入无锁读临界区时声明自己处于当前 epoch；节点 retired 时记录 retire epoch；当所有活动线程都离开旧 epoch 后，旧节点可以回收。它的 fast path 通常比 hazard pointer 轻，但一个线程如果长时间停在临界区或被暂停，可能阻止旧 epoch 回收，导致 retired list 增长。

\begin{lstlisting}[numbers=none,caption={epoch 回收的状态机}]
ThreadOutside
  -> enter critical section, publish local epoch
ThreadActive(epoch)
  -> read lock-free structure
  -> leave, clear active flag
ThreadOutside

NodeLinked
  -> removed by successful CAS
Retired(epoch)
  -> all active threads are newer or inactive
Reclaimed
\end{lstlisting}

Epoch 适合短读临界区和受控线程集合。若线程可能在临界区里做 I/O、等待 future、阻塞系统调用或长期暂停，epoch 回收会失去边界。工程报告要记录最大 retired list、最长 active epoch、线程退出清理和压力下内存峰值。

\section{RCU：读多写少的另一种选择}

RCU 的核心是让读者几乎不阻塞，写者发布新版本，并等待 grace period 后回收旧版本。它适合读多写少、读路径极短、更新可以复制或延迟发布的场景。RCU 不是通用队列替代品。若更新频繁、对象很大或读者可能长期停留，写者延迟和内存占用会增加。

教学项目中，RCU 可用于配置表、路由表或模型版本指针：请求读取当前版本，更新线程构造新版本后原子发布，旧版本等无读者后释放。它不适合把每个 token 的 KV cache 更新都做成 RCU；那会把高频写路径变成版本垃圾制造器。

\section{进展证明要包含 allocator 和等待策略}

很多论文算法假设内存分配不会阻塞，节点不会耗尽，线程不会在临界区暂停。工程实现不能忽略这些条件。若 push 在循环中调用 \code{new}，而 allocator 内部拿锁或 mmap，那么整体进展性质已经依赖 blocking allocator。若队列满时自旋等待，系统在下游停顿时会烧 CPU。若关闭路径等待所有线程退出，取消和析构也属于进展证明的一部分。

\begin{lstlisting}[numbers=none,caption={进展证明字段}]
progress_proof:
  public_operations
  shared_atomic_fields
  possible_retry_loops
  allocation_in_fast_path
  full_or_empty_behavior
  blocking_calls_inside_operation
  cancellation_and_close_path
  starvation_risk
  backoff_strategy
\end{lstlisting}

\section{测试矩阵：无锁需要压力和反例}

无锁测试不能只跑一次吞吐。至少要覆盖：

\begin{itemize}
  \item 单线程退化路径：空、满、容量边界、环绕。
  \item 多线程压力：不同生产者/消费者比例、不同线程绑定、不同输入速率。
  \item 生命周期压力：对象池复用、快速 allocate/free、retired list 增长。
  \item 取消和关闭：生产者退出、消费者退出、drain、stop。
  \item sanitizer：TSAN 能发现部分数据竞争，但不能证明线性化。
  \item model/litmus：小容量、小步数穷举或随机调度重放。
  \item 指标：CAS 失败次数、平均重试、p99 重试、retired list 峰值、CPU 占用。
\end{itemize}

\section{验收标准}

一个无锁结构进入默认路径前，必须提交：

\begin{enumerate}
  \item 公开接口语义和失败语义。
  \item 每个操作的线性化点。
  \item 每个普通数据发布的 release/acquire 证明。
  \item ABA 和内存回收策略。
  \item 进展保证和不满足保证的条件。
  \item 压力测试、sanitizer、模型测试或故障重放。
  \item 与 mutex/condition variable 版本的吞吐、CPU 占用和尾延迟对比。
\end{enumerate}

\begin{keyidea}
无锁结构的门槛不在 CAS 语法，而在证明。不能写出线性化点、回收安全和进展边界的无锁代码，不应进入生产路径。
\end{keyidea}
